%-*-latex-*-
%-*-latex-*-
\newcommand\COURSE{ciss240}
\newcommand\ASSESSMENT{q1103}
\newcommand\ASSESSMENTTYPE{Quiz}
\newcommand\POINTS{\textwhite{xxx/xxx}}

\input{myquizpreamble}
\input{yliow}
\input{\COURSE}
%-*-latex-*-

%-*-latex-*-

%-*-latex-*-
\renewcommand\TITLE{\ASSESSMENTTYPE \ \ASSESSMENT}

\renewcommand\EMAIL{}
\renewcommand\AUTHOR{eagraham1@cougars.ccis.edu}

\textwidth=6in
\begin{document}
\topmatterthree

This is a closed-book, no compiler, 5 minute quiz.

%------------------------------------------------------------------------------
\nextq
Write a program that prints solutions to
the $x^3 + 3y^3 = nz^3$ where
$n$ is an integer input from the user and
$x,y,z$ are integers in the interval $[-10, 10]$.
After the solutions are printed, the number of solutions is printed.
Here's a test case:
\begin{Verbatim}[frame=single,fontsize=\footnotesize,commandchars=\\\{\}]
\userinput{5}
-10 5 -5
-8 4 -4
-6 3 -3
-4 2 -2
-2 1 -1
0 0 0
2 -1 1
4 -2 2
6 -3 3
8 -4 4
10 -5 5
11
\end{Verbatim}
\ANSWER
\begin{answercode}
#include <iostream>
#include <cmath>

int main()
{
    int n;
    std::cin >> n;

    int count = 0;

    for (int x = -10; x <= 10; ++x)
    {
        for (int y = -10; y <= 10; ++y)
        {
            for (int z = -10; z <= 10; ++z)
            {
                if (pow(x, 3) + 3 * pow(y, 3) == n * pow(z, 3))
                {
                    std::cout << x << " " << y << " " << z << '\n'; 
                    count++;
                }
            }
        }
    }

    std::cout << count << '\n';

    return 0;
}

\end{answercode}

%------------------------------------------------------------------------------
\nextq
Write an ASCII art program that does the following.
3 test cases are shown.
\begin{Verbatim}[frame=single,fontsize=\footnotesize,commandchars=\\\{\}]
\userinput{3}
  *
 *
*
\end{Verbatim}
\begin{Verbatim}[frame=single,fontsize=\footnotesize,commandchars=\\\{\}]
\userinput{4}
   *
  *
 *
*
\end{Verbatim}
\begin{Verbatim}[frame=single,fontsize=\footnotesize,commandchars=\\\{\}]
\userinput{5}
    *
   *
  *
 *
*
\end{Verbatim}
\ANSWER
\begin{answercode}
#include <iostream>

int main()
{
    int n;
    std::cin >> n;

    for (int i = 0; i < n; ++i)
    {
        for (int j = 0; j < n - i - 1; ++j)
        {
            std::cout << " ";
        }
        std::cout << "*";
        std::cout << '\n';
    }

    return 0;
}
\end{answercode}

%------------------------------------------------------------------------------
%-*-latex-*-
\newpage
\input{instructions}
\end{document}

